\documentclass[10pt,a4paper]{amsart}
\usepackage[margin=19mm]{geometry}
\usepackage{amsmath,amssymb,mathtools}
\usepackage[scr=rsfs,cal=euler]{mathalfa}
\usepackage{xfrac}
\usepackage[unicode,hidelinks,bookmarks=false]{hyperref}
\newcommand{\ie}{i.e.\ }
\newcommand{\cf}{cf.\ }
\newcommand{\resp}{respectively\ }
\newcommand{\Subsection}{\subsection}
\providecommand{\nobreakdash}{\nobreak\hbox{-}}
\setlength{\emergencystretch}{2em}
\multlinegap0pt
\usepackage{float}
\usepackage{amssymb}
\usepackage{xcolor}
\usepackage{algorithm}\usepackage{algpseudocode}
\newtheorem{proposition}{Proposition}
\newcommand{\bx}{\mathbf{x}}
\newcommand{\mA}{\mathcal{A}}
\newcommand{\mABX}{\mathcal{A} \in \mathcal{B}(\mathcal{X})}
\newcommand{\mBX}{ \mathcal{B}(\mathcal{X})}
\newcommand{\mM}{\mathcal{M}}
\newcommand{\mS}{\mathcal{S}}
\newcommand{\mT}{\mathcal{T}}
\newcommand{\mX}{\mathcal{X}}
\newcommand{\ov}{\overline{v}}
\title{\bfseries {\Large Optimizing Treatment Allocation to Maximize the Health of a Population}}
\author{Daniel Adelman, Alba V Olivares-Nadal, Miaolan Xie}
\date{Source: arXiv:2604.07738v1 (CC BY 4.0)}
\begin{document}
\maketitle

\noindent\emph{Source and licence.} \bfseries {\Large Optimizing Treatment Allocation to Maximize the Health of a Population}. Version arXiv:2604.07738v1 (CC BY 4.0), distributed by arXiv under the Creative Commons Attribution 4.0 International licence, http://creativecommons.org/licenses/by/4.0/, as recorded in the arXiv licence metadata for this exact version and shown on the abstract page https://arxiv.org/abs/2604.07738v1. Full licence notice: \texttt{licenses/arxiv-2604.07738-CC-BY-4.0.txt}.\par\smallskip

\noindent\textit{Editorial scope.} Counted unit: the article's proposition lemma:phi (source0.tex lines 495--497). The article's complete original proof of it is delivered with it (source0.tex lines 499--530, one \texttt{proof} environment of 32 lines). No mathematics is written, repaired or re-derived: the statement and the proof are the article's own lines, in the article's own notation, and no retained line is substituted or re-flowed.  Retained uncounted as the source-local context the proof needs: lines 382--385; lines 431--434; lines 477--482.  Nothing was omitted from any retained span, so the package carries no omission record.  No retained line contains a citation, so the wrapper carries no bibliography and no bibliography entry.  No retained line contains a figure environment, an image-inclusion command or drawing code, so no image is delivered and the package declares no asset dependency.  Editorial formatting only, in addition: an AMS \texttt{amsart} preamble that restates the article's own declarations and macros used by the retained lines, namely the environments \texttt{proposition} on separate counters, together with the standard packages those lines need.  The retained environments print in this document as Proposition 1 in order of appearance, whereas the article prints the same items with the numbering of its own full document.  The complete native source of the article as distributed in the arXiv e-print for this exact version is bundled unchanged as \texttt{sources/198111/source0.tex}.
	\begin{align}\label{K*1}\tag{K$^*_1$}
		\ov^*(\nu)= \ \max_{\substack{\varphi \in \Phi}}& \ \mathbb{E}_{\nu}\mathbb{E}_{\varphi} [(z+u) y_1(\bx)]+\mathbb{E}_{\nu}\mathbb{E}_{\varphi} [(1-z-u) y_0(\bx)] + \alpha \ov^{*}(f(\nu,\varphi))\\
		\mbox{s.t.} & \  \mathbb{E}_{\nu}\mathbb{E}_{\varphi} [z+u] \leq m/n, \nonumber
	\end{align}

	\begin{align}
		\ov^{*}(\nu)= \ \max_{\substack{\varphi \in \Phi}} &\  \int_{\mX }  y_1(\bx)   \left\{ d\nu(\bx,1)  +\varphi(1| \bx,0) d\nu(\bx,0)\right\}   + \int_{\mX }  y_0(\bx)    \varphi(0| \bx,0) d\nu(\bx,0)   +\alpha \ov^{*}(f(\nu,\varphi)) \nonumber\\
		\mbox{s.t.}  & \ \   \nu(\mX,1)+\int_{\mX} \varphi(1|\bx,0)d\nu(\bx,0)  \leq m/n. \label{K*2}
	\end{align}

	\begin{align}\label{K*tau} \tag{K$^*_\tau$}
		\ov^*(\eta,\rho)=& \ \max_{\substack{\tau \in \mT(\eta,\rho)}} \ \mathbb{E}_{\rho+\tau}\left[ y_1(\bx) \right]     +\mathbb{E}_{\eta-\tau}\left[  y_0(\bx) \right]   + \alpha \ov^*(f((\eta,\rho),\tau)), \qquad \forall (\eta,\rho)\in \mS_\mM.
        %\\
		%\mbox{s.t.} & \ \   \rho(\mX)+\tau(\mX)  \leq m/n. \nonumber\\
		%& \ \  \tau(\mA)  \leq \eta(\mA) \nonumber
	\end{align}

	\begin{proposition}\label{lemma:phi}
		Fix a state $(\eta,\rho)$; then for every solution $\tau$ of \eqref{K*tau}, there exists a solution $\varphi \in \Phi$ of \eqref{K*1} that is unique $\eta$-almost everywhere.
	\end{proposition}

	\begin{proof}{Proof}
		We start by proving that there exists a measurable function $g: \mathcal{X} \rightarrow [0,1]$ such that
		\begin{equation} \label{radon-nikodym}
			\tau(\mA)=\int_{\mA} g(\bx) d\eta(\bx) \qquad \forall \mABX.
		\end{equation}
		
		The second constraint in \eqref{K*tau} implies that $\tau$ is absolutely continuous with respect to $\eta$, since $\tau(\mA)=0$ for every $\mABX$ such that $\eta(\mA)=0$. Then the Radon-Nikodym theorem states that there exists a measurable function $g: \mathbb{R} \rightarrow [0,\infty)$ such that \eqref{radon-nikodym} holds, and $g$ is unique $\eta$-almost everywhere. 
		
		
		To prove that $g$ maps $\mX$ into the interval $[0,1]$, define $\mA_\epsilon=\{ \bx \in \mX: \ g(\bx)\geq 1+\epsilon \}$ and assume that $\exists \ \epsilon'>0$ such that $\eta(\mA_{\epsilon'})>0$. Since $(\mX,\mBX)$ is a measurable space and $g$ is a measurable function in that space, then $\mA_\epsilon \in \mBX$ for all $\epsilon>0$. Then
		
		$$\tau(\mA_{\epsilon'})=\int_{\mA_{\epsilon'}} g(\bx) d\eta(\bx) \geq (1+ \epsilon') \cdot \eta(\mA_{\epsilon'}) > \eta(\mA_{\epsilon'}) $$
		
		\noindent where the last inequality is strict due to the fact that we assumed that $\eta(\mA_{\epsilon'})>0$. Since this is a contradiction with the second constraint \eqref{K*tau}, we can claim that function $g$ is between 0 and 1 $\eta$-almost everywhere. Define the solution 
		
		$$
		\varphi(u|\bx,z)=\left\{ 
		\begin{array}{ll}
			g(\bx) & \quad \mbox{ if } u=1,z=0\\
			1-g(\bx) & \quad \mbox{ if } u=0,z=0\\
			0 & \quad \mbox{ if } u=1,z=1\\
			1 & \quad \mbox{ if } u=0,z=1
		\end{array}
		\right.
		$$
		
		
		Let us show that $\varphi$ is a stochastic kernel. First,  for any fixed action $u \in \{0,1\}$, the function $\varphi(u|\cdot,\cdot)$ is measurable because it is either a constant , $g$ or $1-g$, which are measurable on $\mX$. Second for any fixed state $(\bx,z) \in \mX \times \{0,1\}$, we need to prove that the function $\varphi(\cdot|\bx,z)$ is a probability measure. It is obvious that $\varphi(\cdot|\bx,z) \in [0,1]$ because $g$ itself is. Finally, $\varphi$ is a solution to \eqref{K*1} because $\tau$ verifies the first constraint in \eqref{K*tau}, so $\varphi$ verifies the constraint in \eqref{K*2}, which is an equivalent formulation of \eqref{K*1}.
		
		
		Now we prove that this solution is unique $\eta$-a.e. If $\exists \ \epsilon'>0$ such that $\mA_{\epsilon'} \neq \emptyset$ then we could construct a function $g': \mX \rightarrow [0,1]$ such that $g'(\bx)\in [0,1]$ for all $\bx \in \mA_{\epsilon'}$ for all $\epsilon' >0$, $g'$ being equal to $g$ $\eta$ -almost everywhere and fulfilling \eqref{radon-nikodym}. \hfill $\square$
	\end{proof}

\end{document}
